\clearpage
\section{Locked evidence audit and bounded next experiment}
\subsection{The corrected benchmark is a negative result}
The committed \texttt{results/s669\_group\_holdout.json} contains 603 training variants and 66 test variants, grouped by PDB accession. On the 66 test rows the GNN MAE is 1.4219 kcal/mol and Spearman correlation is 0.4024. Stored GeoDDG-Seq predictions on the same rows have MAE 1.3101 and Spearman 0.4030. This is not a stability-prediction win. The earlier 420-row subset came from a numbering error and cannot be used as a competing headline result. A single random group split also cannot quantify the variability of training across all plausible protein splits. The model is a useful test of a simple graph representation and a source of error analyses, but not a validated redesign tool.

The committed \texttt{results/group\_sensitivity.json} gives GNN-minus-GeoDDG-Seq absolute-error differences of +0.1118 kcal/mol when each mutation is weighted equally and -0.0413 kcal/mol when each protein is weighted equally. Their signs disagree. The protein-equal bootstrap interval [-0.4337, +0.3516] crosses zero. Both the row and protein summaries are descriptive views of the same fixed test set, not independent confirmation. The sparse-group analysis and leave-one-PDB-out diagnostics in this manuscript show that the sign of a selected aggregate depends on which proteins receive weight. No post-hoc subgroup should be promoted into a prospective rescue claim.

\subsection{What the disease-target material can establish}
The target structures and cited suppressor studies supply case context, not proof that this graph model found a functional rescue in TP53, SOD1, or PTEN. A predicted stability change alone does not measure folding in the mutant background, proteostasis, transcriptional function, or disease phenotype. Wild-type-background single mutants and double mutants on a disease background answer different biochemical questions. The current analysis lacks prospective double-mutant measurements. [PENDING: a held-out mutant-background rescue panel with measured stability and function, predeclared before model ranking.]

\subsection{Data use and provenance}
The source audit distinguishes complete S669 graph construction from the partial FireProtDB export. A direct number match and residue check on 1,470 FireProtDB rows does not make those rows harmonized with S669 assay conditions. The source is a partial snapshot, not a census. In particular, the 157 unique mutation-data PDB accessions and six separately studied disease structures are structure-use counts; they are not 163 independent experimental cohorts. [PENDING: a complete, independently sourced and assay-harmonized external validation set.] The result tables should retain this denominator distinction throughout.

\subsection{Falsifiable decision gate}
A follow-up would freeze a split by protein family before training, archive source identifiers and mutation-to-structure mappings, and compare paired predictions from established tools on the exact same assay rows. It would then test rescue nominations on disease-mutant backgrounds rather than selecting examples after observing favorable scores. The current report does not contain that experiment. Its defensible result is the corrected mapping audit and a negative generalization comparison, with uncertainty driven in part by a small and uneven protein-held-out test set.
